% 第9章 デリバティブの価格付け(1)(資料9)
\chapter{デリバティブの価格付け (1):無裁定価格とモデルフリーな価格式}

本章から,株式デリバティブの価格付け (pricing) を離散時間モデルで学ぶ.
価格付けの基礎となるのは\textbf{裁定機会}と\textbf{無裁定価格}の考え方,そして
\textbf{複製}という手法である.本章ではこれらの概念を導入し,モデルに依存しない
(モデルフリーな)価格式の例として,先渡価格と金利スワップのスワップレートを
導出する.

\section{裁定機会と無裁定価格}

\subsection{金融商品の価格付けとは}

\textbf{金融商品}とは,将来のキャッシュフロー列の交換に関する契約である.
その\textbf{価格付け}とは,将来に発生する(不確実な)キャッシュフロー列の
現在価値を決定することである.

なお,将来のキャッシュフロー列に不確実性がなければ,評価は簡単である.
リスクのない割引債のポートフォリオとみなせばよい.すなわち,各時点の
キャッシュフローに分解し,それぞれのキャッシュフローを割引債とみなして
割引債価格の合計値を求めればよい(第3章で説明済).問題は,キャッシュフローが
\emph{不確実}な場合である.

\subsection{裁定機会}

以下では市場の完全性(完全市場)を仮定する.

\begin{definition}[完全市場(再掲)]
\begin{enumerate}
\item[①] 個々の投資家はプライステイカーである(個々の投資家の取引は市場価格に
影響を与えない).
\item[②] 証券は無限分割可能(任意の量の取引が可能)である.
\item[③] 取引費用や税金は存在しない.
\item[④] 投資家は自由に空売りができる(負の保有量も可能).
\end{enumerate}
ここで\textbf{空売り}とは,もともと保有していない証券を(借りてきて)売ることで
ある.先に代金を受け取り,後日(満期日に)買い戻す.資産の保有量としては
「マイナス」(負)とみなされる.
\end{definition}

\begin{definition}[裁定機会]
\textbf{裁定機会}とは,コストゼロで,損失を被ることなく,利益をあげることが
できる(または,利益をあげる\emph{可能性のある})投資機会のことである.
「必ず利益をあげられる」でなくてよい.利益を得る可能性さえあればよい.
\end{definition}

もしも市場に裁定機会が存在したら,どうなるか.多くの人がその裁定機会を使って
儲けようとして取引を行う.すると需給関係が変化し,裁定機会にある資産の市場価格が
変わる.やがて裁定機会は消滅する.

\begin{example}[裁定機会の例]
ある満期 $T$ で,キャッシュフロー $X$ が得られる証券と,キャッシュフロー $Y$ が
得られる証券があるとする($X = $ 林檎1個,$Y = $ 蜜柑1個,と思ってもよい).
それぞれの価格(現在価値)を $\pi(X)$, $\pi(Y)$ とする.また,同じ満期 $T$ で
キャッシュフロー $X+Y$ が得られる証券があり,その価格を $\pi(X+Y)$ とする.
\begin{itemize}
\item $\pi(X) + \pi(Y) < \pi(X+Y)$ ならば,裁定機会が存在する.この場合の裁定
機会は,「$X$ と $Y$ を別々に買い,$X+Y$ の証券を空売りする」ことである.
現時点で差額 $\pi(X+Y) - \pi(X) - \pi(Y) > 0$ を受け取り,満期のキャッシュ
フローは完全に相殺される.
\item $\pi(X) + \pi(Y) > \pi(X+Y)$ ならば,やはり裁定機会が存在する.この場合は
逆に「$X+Y$ の証券を買い,$X$ と $Y$ を別々に空売りする」ことである.
\end{itemize}
当初は裁定機会があったとしても,十分に取引が進めば,
\[
\pi(X) + \pi(Y) = \pi(X+Y)
\]
という状態になると考えられる.これが「裁定機会が存在しない」状態である.
\end{example}

\begin{keypoint}
効率的な市場では,このような(裁定機会が消えていく)変化が一瞬で進むと考えて,
近似的に市場には\textbf{「裁定機会が存在しない」}とみなす.
\end{keypoint}

なお,スーパーマーケットなどで見られるような「リンゴ1個で100円,10個で800円」
といった取引は考えないこと.「儲かる取引ならば何でもやってやるぞ」といった
人達が集った市場だと思うことが大切である.

\subsection{無裁定価格}

では,適正な価格とは何か.

\begin{definition}[無裁定価格]
\textbf{無裁定価格}とは,裁定機会が存在しないように付けられた価格のことである.
これはファイナンスにおける\textbf{一物一価の法則}である.
\end{definition}

$\pi(X)+\pi(Y) = \pi(X+Y)$ が意味することは,「\textbf{将来キャッシュフローが
等しい2つの証券の価格は等しい}」ということである.価格がわからないデリバティブが
あるとき,その将来キャッシュフローと完全に等しい将来キャッシュフローを持つ
ポートフォリオがあれば,デリバティブの価格はそのポートフォリオの価格に等しい
はずである.このことを,以下でもう少しきちんと表現する.

\subsection{無裁定価格の線形性}

金融商品の無裁定価格は\textbf{線形性}を持つ.金融商品 $X$, $Y$ の無裁定価格を
$\pi(X)$, $\pi(Y)$ とすると,
\begin{equation}
\pi(aX + bY) = a\,\pi(X) + b\,\pi(Y),\qquad a, b \text{ は定数}
\end{equation}
が成り立つ.また,複数時点 $T_i$($i=1,2,\dots,n$)のキャッシュフローが
$C(T_i)$ である金融商品の価格 $\pi$ は,
\begin{equation}
\pi\bigl(C(T_1), C(T_2),\dots,C(T_n)\bigr) = \sum_{i=1}^{n}\pi\bigl(C(T_i)\bigr)
\end{equation}
となる.将来CFが確定しているケースでは,これは当然の話である(第3章).将来CFが
確定して\emph{いない}ケースでも,これは成り立つと考える.

\subsection{複製と複製ポートフォリオ}

\begin{definition}[複製]
金融資産 $X$ のキャッシュフローと全く同じキャッシュフローを生成するポート
フォリオ $Y$ が存在するとき,$Y$ は $X$ を\textbf{複製する}という.このとき,
$Y$ を $X$ の\textbf{複製ポートフォリオ}という.
\end{definition}

\begin{keypoint}[無裁定による価格付けの基本原理]
金融資産 $X$ の無裁定価格は,その複製ポートフォリオ $Y$ の価格(現在価値)に
等しい.

これが無裁定による価格付けの最も基本的な考え方である.金融商品 $X$ の価格を
知るには,(既に価格がわかっている資産で構成される)複製ポートフォリオ $Y$ を
何らかの方法で求めればよい.金融商品 $X$ の価格は,複製ポートフォリオ $Y$ の
価格に等しい.(正確には,もう少し詳しい条件(自己充足性,後述)が必要である.)
\end{keypoint}

\paragraph{「複製」の意味.}
複製とは,ある金融商品(デリバティブ)の価格の推移と,(価格が既知の)資産から
なる複製ポートフォリオの価格の推移が,\textbf{常に一致}することを意味する.
複製ポートフォリオの資産構成は時間とともに変化していく(動的に組み替えていく)
が,その価値は原資産価格のどのような動き方(\textbf{サンプルパス})に対しても,
デリバティブの価値と一致し続け,満期にはペイオフに等しい価値に到達する.
「すべてのサンプルパスにおいて」成立することがポイントである.

\subsection{自己充足性}

この議論のために必要なポートフォリオの性質が,\textbf{自己充足(自己資金調達)
性}である.

\begin{definition}[自己充足性]
「ポートフォリオ」と「(ポートフォリオの)外部」の間で資金の流出入がないとき,
ポートフォリオは\textbf{自己充足}(あるいは自己資金調達的,self-financing)で
あるという.
\end{definition}

\begin{itemize}
\item 何かの証券を購入したければ,既に保有している資産を売却して,その代金を
作らなければならない.そして,手数料も税金も存在しない.
\item つまり,ポートフォリオをリバランス(組み替え)してもポートフォリオ全体の
価値は変わらない.
\end{itemize}

2資産(価格 $(B_t, S_t)$,保有量 $(b_t,\Theta_t)$)のケースでは,時間の流れは
「価格変化 $\to$ 組み替え $\to$ 価格変化 $\to$ …」の繰り返しであり,組み替えの
瞬間にポートフォリオの価値が変わらないこと(自己充足性)が要求される.

価格付けの原理をファイナンスの用語を使って言い直すと:
\begin{keypoint}
デリバティブ $X$ の価格は,その\textbf{自己充足な複製ポートフォリオ} $Y$ の
価格に等しい.
\end{keypoint}

自己充足な複製ポートフォリオであれば,
\begin{itemize}
\item 自分でやりくりする($=$保有資産を売買する)だけで,
\item 外部から資金を受け入れることなく,また,
\item 外部に資金を放出することなく,
\end{itemize}
デリバティブの満期時点のペイオフを再現できる.それも,\emph{どのようなサンプル
パスのときも}再現できる.だからこそ,「複製ポートフォリオの現在価値は
デリバティブの現在価値に等しい」と言えるのである.

以下では,複製ポートフォリオはすべて自己充足的として話を進める.では,どうやって
その複製ポートフォリオを求めるか.以下,具体例で見ていく.

\section{先渡価格}

本節と次節では,\textbf{モデルフリーな(モデルに依存しない)価格式}を導出する.
原資産価格の確率モデルを仮定しなくても,無裁定の考え方だけで価格が決まる例で
ある.

\subsection{先渡契約(再確認)}

\textbf{先渡(フォワード)契約}とは,ある資産(原資産)を,将来の決められた時点
(満期)に,現時点で決められた価格(先渡価格)で売買する契約である(第8章参照.
例:A社の株式1株(現在90万円)を1年後に95万円で売買).

先渡では契約時点(現時点)で金銭のやり取りをしないから,\textbf{この契約の
価値は現時点でゼロ}であり,現在価値がゼロになるように先渡価格が決まる.そして
満期には,\emph{必ず}株式と代金を交換しなければならない.

\subsection{先渡契約の複製ポートフォリオ}

現在時刻を $t$,時刻 $t$ における証券価格を $S(t)$ とする.

\paragraph{先渡契約のキャッシュフロー.}
ロング・ポジションを考える.満期 $T$ に,先渡価格 $K$ を支払い,価格 $S(T)$ の
証券を受け取る.つまり,満期 $T$ で受け取る将来キャッシュフロー(将来価値)は
\begin{equation}
S(T) - K
\end{equation}
である.これは $S(T)$ の1次関数であり,$S(T)=K$ を境に正にも負にもなる直線の
ペイオフである.複製ポートフォリオはこのキャッシュフローの価値を常に実現しな
ければならない.

\paragraph{具体的な複製ポートフォリオ.}
次のポートフォリオを考える.
\begin{itemize}
\item 現時点 $t$ で,証券1単位(価格 $S(t)$)を購入し,時刻 $T$ まで保有する.
\item 現時点 $t$ で,満期 $T$ の割引債(価格 $v(t,T)$)を額面 $K$ だけ空売り
する.
\end{itemize}
(割引債の額面は1円とする.理論ではこの設定をよく使う.)
このポートフォリオの満期 $T$ における価値は
\[
S(T) - K \times 1 = S(T) - K
\]
であり,先渡契約のロング・ポジションを複製している.

\subsection{先渡価格の導出}

先渡契約では,現時点でキャッシュフローは発生しない.よって\textbf{先渡契約の
現在価値($=$価格)はゼロ}である.したがって,複製ポートフォリオの現在価値
(価格)もゼロでなければならない.このことを数式で表すと,時刻 $t$($t<T$)に
おいて
\begin{equation}
S(t) - K\,v(t,T) = 0.
\end{equation}
左辺は複製ポートフォリオの現在価値(価格)を与える式,右辺は先渡契約の現在価値
(価格)である.よって,合理的な\textbf{先渡価格}は
\begin{equation}
K = \frac{S(t)}{v(t,T)}
\label{eq:forward-price}
\end{equation}
である.通常(金利がプラス)ならば $v(t,T) < 1$ なので,$K > S(t)$ となる
(先渡価格は現在の株価より高い).

\section{金利スワップの価格付けとスワップレート}

これもモデルフリーな価格式の例である.

\subsection{金利スワップ(再確認)}

スワップ取引とは,将来のキャッシュフロー列を別のキャッシュフロー列と交換する
取引である(金利スワップ,通貨スワップ,ベーシス・スワップ等.第8章参照).
典型例の金利スワップでは,ある共通の元本(\textbf{想定元本})に対する
「固定金利払」と「変動金利払」を交換する.交換は金利払の分のみで,想定元本は
交換しない.変動払のキャッシュフローは,各時点ごとに決まる金利に応じて決まる
ものとする.金利スワップは,変動金利による借入・借換と深い関係がある.

\subsection{変動金利による借入}

1円を変動金利 $L(t)$ で期間 $\delta$ で借り入れる.そのキャッシュフローは,
時刻 $t$ に1円の受取,時刻 $T = t+\delta$ に元本1円と利息 $\delta L(t)$ の支払
である.ポイントは,\textbf{利率 $L(t)$ は時刻 $t$ に決まる}ことである.

このような借入(貸付)において,受取と支払のキャッシュフローの現在価値は等しいと
考える(だから異時点間でキャッシュの交換を行える).受取側のキャッシュフローの
現在(時刻 $t$)における価値は1円である.なので,\textbf{支払側のキャッシュ
フロー(時刻 $t+\delta$ の $1+\delta L(t)$)の現在(時刻 $t$)における価値も
1円}である.

\subsection{変動払と変動金利による借換}

\paragraph{変動払の複製.}
期間 $\delta$ ごとに変動金利で1円を借換えると,そのキャッシュフローは,
各時点 $t_i$ で「元本1円の返済+利息 $\delta L(t_{i-1})$ の支払,同時に新たに
1円の借入」の繰り返しになる.元本部分の受取・支払は各時点で相殺されるので,
これは次のキャッシュフローと同じである:
\[
\text{時刻 } t_0 \text{ に1円の受取},\quad
\text{時刻 } t_i \text{ に } \delta L(t_{i-1}) \text{ の支払}\ (i=1,\dots,4),
\quad
\text{時刻 } t_4 \text{ に元本1円の支払}.
\]

このキャッシュフローのうち,利息部分(時刻 $t_1,\dots,t_4$ の
$\delta L(t_0),\dots,\delta L(t_3)$)は\textbf{金利スワップの変動払と同じ}で
ある.受取と支払のキャッシュフローの現在価値は等しいので,
\[
\text{(時刻 } t_0 \text{ の1円受取の現在価値)}
= \text{(利息CFの現在価値)} + \text{(時刻 } t_4 \text{ の元本1円支払の現在価値)}
\]
が成り立つ.$v(t,T)$ を時刻 $t$ における満期 $T$ の割引債価格とすると,
このことから
\begin{equation}
\text{(変動払CFの現在価値)} = v(t,t_0) - v(t,t_4)
\end{equation}
が得られる.

\subsection{金利スワップの評価}

\paragraph{記号の設定.}
固定金利 $S$ と変動金利 $L(t)$ の交換を考える.時刻 $t_0$ 開始,満期 $t_m$,
利払日 $t_1,t_2,\dots,t_m$(間隔 $\delta$:一定),想定元本 $A=1$ とする.
\begin{itemize}
\item \textbf{固定払}:時刻 $t_1,t_2,\dots,t_m$ に $\delta S$ を支払う.
\item \textbf{変動払}:時刻 $t_i$($i=1,2,\dots,m$)に $\delta L(t_{i-1})$ を
支払う.
\end{itemize}

\paragraph{変動払の価値.}
前項のロジックにより,時刻 $t_i$($i=1,\dots,m$)に $\delta L(t_{i-1})$ を
支払う変動払の(時刻 $t$ における)価値は
\begin{equation}
v(t,t_0) - v(t,t_m)
\label{eq:float-leg}
\end{equation}
である.

\paragraph{固定払の価値.}
固定払の複製は,時刻 $t_i$($i=1,\dots,m$)を満期とする割引債(額面
$\delta S$)のポートフォリオである.その価値は
\begin{equation}
\sum_{i=1}^{m}\delta S \times v(t,t_i)
= \delta S\sum_{i=1}^{m} v(t,t_i).
\label{eq:fixed-leg}
\end{equation}

\paragraph{金利スワップの価値.}
時刻 $t$ における,時刻 $t_0$ 開始・満期 $t_m$ の金利スワップの価値は,
固定受変動払側からみて(立場によって価格は符号が逆になる),
\begin{equation}
\delta S\sum_{i=1}^{m} v(t,t_i) - \bigl[v(t,t_0) - v(t,t_m)\bigr].
\end{equation}

\paragraph{スワップレート.}
契約時点の金利スワップではキャッシュの交換はないので,
\[
\text{(スワップの固定払の価値)} = \text{(スワップの変動払の価値)},
\]
すなわち
\begin{equation}
\delta S\sum_{i=1}^{m}v(t,t_i) = v(t,t_0) - v(t,t_m)
\end{equation}
が成り立つ.よって,時刻 $t$ における時刻 $t_0$ 開始,満期 $t_m$ の
\textbf{スワップレート}は
\begin{equation}
S(t, t_0, t_m; \delta)
= \frac{v(t,t_0) - v(t,t_m)}{\delta\sum_{i=1}^{m}v(t,t_i)}
\label{eq:swap-rate}
\end{equation}
である.

\begin{keypoint}
金利期間構造($=$割引債価格の期間構造)が決まると,それに応じて,開始時点
$t_0$,満期 $t_m$,利払日の間隔 $\delta$ のスワップレートが一意に決まる.
$\delta$ を\textbf{テナー} (tenor) という.

なお,実際のマーケットでは上記と\emph{逆}で,スワップレートの期間構造の情報
から,割引債価格の期間構造を得る.
\end{keypoint}

\begin{remark}[金融危機以後の注意]
以上の話は,金融危機以前の話である.金融危機以後,(特に金利)デリバティブの
評価は複雑化した.
\begin{itemize}
\item これまではテナーが異なっても同じ金利期間構造($=$割引債価格の期間構造)を
もとにスワップレートが決まったが,この関係が崩れた.わずかな信用リスクの違いを
評価するようになったためである.
\item その結果,実務ではここまでの話では通用しなくなった.
\item さらに,無リスク変動金利はLIBORからON (overnight) 金利になった.
金融は日々変化する.
\end{itemize}
とはいえ,価格付けの基礎を理解するには,前述の金利スワップの考え方,スワップ
レートの算出ロジックは今も重要である.
\end{remark}

% ---------------- 演習問題 ----------------
\section{演習問題}

\begin{exercise}[問題1]
期間 $t$ 年のゼロレートを $R(t)$ とする.$R(1)=3.0\%$, $R(2)=5.0\%$,
$R(3)=6.0\%$, $R(4)=6.5\%$, $R(5)=7.0\%$ として,満期1, 2, 3, 4, 5年の
スワップレートを小数第2位まで求めなさい.ただし,テナーは1年とし,連続時間
モデルで考えること.
\end{exercise}

\begin{answerbox}
\textbf{考え方.}スワップレートの式\eqref{eq:swap-rate}を,現時点 $t=0$ 開始
($t_0 = t = 0$,このとき $v(0,0)=1$),テナー $\delta=1$ 年で使う:
\[
S_m = \frac{1 - v(0,m)}{\sum_{i=1}^{m}v(0,i)}.
\]
割引債価格は連続時間モデル(連続複利)なので $v(0,t) = e^{-R(t)\,t}$ で計算
する.

\medskip
\textbf{計算.}
まず割引債価格を求める:
\[
v(0,1) = e^{-0.03} = 0.970446,\quad
v(0,2) = e^{-0.10} = 0.904837,\quad
v(0,3) = e^{-0.18} = 0.835270,
\]
\[
v(0,4) = e^{-0.26} = 0.771052,\quad
v(0,5) = e^{-0.35} = 0.704688.
\]
分母の累積和とスワップレートをまとめると次表のとおりである.
\begin{center}
\begin{tabular}{ccccc}
\toprule
$t$ (yrs) & $R(t)$ & $v(0,t)$ & 分母 $\sum_{i=1}^{t}v(0,i)$ & Swap Rate\\
\midrule
0 & --- & 1.000000 & --- & ---\\
1 & 3.00\% & 0.970446 & 0.970446 & 3.05\%\\
2 & 5.00\% & 0.904837 & 1.875283 & 5.07\%\\
3 & 6.00\% & 0.835270 & 2.710553 & 6.08\%\\
4 & 6.50\% & 0.771052 & 3.481605 & 6.58\%\\
5 & 7.00\% & 0.704688 & 4.186293 & 7.05\%\\
\bottomrule
\end{tabular}
\end{center}
例えば満期2年では
\[
S_2 = \frac{1 - 0.904837}{0.970446 + 0.904837} = \frac{0.095163}{1.875283}
= 5.07\%
\]
である.

\medskip
\textbf{ポイント.}スワップレートはゼロレートと近い値になるが,一致はしない.
スワップレートは「割引債価格で加重平均された金利水準」のような量であり,
ゼロレートカーブが右上がりのとき,満期 $m$ のスワップレートは満期 $m$ の
ゼロレートよりわずかに高くなっている(手前の低い金利の割引債の重みが大きい
ため,分母が大きくなりにくいことによる).
\end{answerbox}
